\section{附注}\label{sec:notes-on-geometric-mechanics}

\subsection{第 \ref{sec:mechanical-systems} 节}\label{subsec:notes-mechanical-systems}

从第~\ref{sec:mechanical-systems}~节的习题开始, 本章将用到一些关于 $\mathbb{R}^n$ 中向量场的不动点稳定性的定义和事实 (更多细节可参见 [\ref{bib:ch05-Arn92}, \ref{bib:GH02}]). 为研究非线性系统
\begin{equation}\label{eq:nonlinear-system-at-fixed-point}
  \begin{cases}
    \dot x=f(x)\\
    x(0)=x_0
  \end{cases}
  \qquad(x\in\mathbb{R}^n),
\end{equation}
通常先求 $f$ 的零点, 称为\textbf{不动点}, \textbf{平衡点}或\textbf{定常解}. 若对 $\overline x$ 的每个邻域 $U$, 都存在 $\overline x$ 的另一个 (可能更小的) 邻域 $V$, 使只要 $x_0\in V$, 在解有定义的每个 $t>0$ 都有 $x(t)\in U$, 则称不动点 $\overline x$ \textbf{稳定}. 在许多情形下, 可以通过在 $\overline x$ 处将 \eqref{eq:nonlinear-system-at-fixed-point} 线性化, 并分析所得的 (线性) 系统
\begin{equation}\label{eq:linearized-system-at-fixed-point}
  \begin{cases}
    \dot\xi=A\xi\\
    \xi(0)=\xi_0
  \end{cases}
  \qquad(\xi\in\mathbb{R}^n)
\end{equation}
来研究 $\overline x$ 附近解的行为, 其中 $A:=(\mathrm{d}f)_{\overline x}$. 此线性系统有全局解
\[
  \xi(\xi_0,t)=e^{tA}\xi_0,
\]
其中可将 $e^{tA}$ 看作由 $\mathbb{R}^n$ 到 $\mathbb{R}^n$ 的映射, 它定义了向量场 $A\xi$ 的流. 若将 $A$ 化为 Jordan 标准形, 则显然此流具有如下不变子空间:
\begin{align*}
  E^s&:=\operatorname{span}\{v_1,\ldots,v_{n_s}\}
       &&\text{(稳定子空间)};\\
  E^u&:=\operatorname{span}\{u_1,\ldots,u_{n_u}\}
       &&\text{(不稳定子空间)};\\
  E^c&:=\operatorname{span}\{w_1,\ldots,w_{n_c}\}
       &&\text{(中心子空间)},
\end{align*}
其中 $v_1,\ldots,v_{n_s}$ 是对应于实部为负的特征值的 $n_s$ 个广义特征向量, $u_1,\ldots,u_{n_u}$ 是对应于实部为正的特征值的 $n_u$ 个广义特征向量, 而 $w_1,\ldots,w_{n_s}$ 是对应于实部为零的特征值的 $n_c$ 个广义特征向量. 若 $E^c=\varnothing$, 则称 $\overline x$ 为 $f$ 的\textbf{双曲}或\textbf{非退化}不动点. 在此情形下, \textbf{Hartman–Grobman 定理}告诉我们, 存在一个从 $\mathbb{R}^n$ 中 $\overline x$ 的某邻域到 $\mathbb{R}^n$ 中 $0$ 的某邻域的同胚, 将 \eqref{eq:nonlinear-system-at-fixed-point} 的非线性流的轨道映为 \eqref{eq:linearized-system-at-fixed-point} 的线性流 $e^{tA}$ 的轨道. 此时 $\overline x$ 附近解的渐近行为, 从而其稳定性类型, 由 $A$ 的特征值 $\lambda$ 决定.

\subsection{文献附注}\label{subsec:geometric-mechanics-literature}

本章材料紧随 \RGcite{Oli02}, \RGcite{Arn97}. 当然, 还有很多优秀的力学著作, 既有传统力学 \RGcite{GPS02}, 也有几何力学 [\ref{bib:AM78}, \ref{bib:MR99}]. 非完整系统 (包括控制理论) 在 [\ref{bib:Blo03}, \ref{bib:BL05}] 中有更详细的讨论. 关于完全可积系统的更多资料, 可参见 [\ref{bib:CB97}, \ref{bib:Aud96}].
